% TODO make hw stylesheet

\documentclass[10pt]{article}
\usepackage[letterpaper,top=1in,left=1in,right=1in]{geometry}
\usepackage[dvipsnames,svgnames]{xcolor}
\usepackage[shortlabels]{enumitem}
\usepackage[framemethod=TikZ]{mdframed}
\usepackage{amsmath,amssymb,amsthm}
\usepackage[colorlinks]{hyperref}
\usepackage{multicol}
\usepackage{tabularx}
\usepackage{bbm}
\usepackage{tikz-cd}

\hypersetup{
    colorlinks=true,
    linkcolor=blue,
    filecolor=magenta,
    urlcolor=magenta,
    pdftitle={MAT 535 HW}
}

\usepackage{mathtools}
\usepackage{thmtools}
\usepackage[textsize=footnotesize]{todonotes}
\usepackage{titling}
\usepackage{cancel}

\reversemarginpar
\mdfdefinestyle{mdbluebox}{roundcorner=10pt,innerbottommargin=9pt,
linecolor=blue,backgroundcolor=TealBlue!5,}
\declaretheoremstyle[headfont=\sffamily\bfseries\color{MidnightBlue},
mdframed={style=mdbluebox},]{thmbluebox}

\mdfdefinestyle{mdbloobox}{frametitlefont=\bfseries,innerbottommargin=8pt,
nobreak=true,backgroundcolor=TealBlue!5,linecolor=blue,}
\declaretheoremstyle[headfont=\bfseries\color{MidnightBlue},
mdframed={style=mdbloobox},headpunct={\\[3pt]},postheadspace=0pt,]{thmbloobox}

\mdfdefinestyle{mdredbox}{frametitlefont=\bfseries,innerbottommargin=8pt,
nobreak=true,backgroundcolor=Salmon!5,linecolor=RawSienna,}
\declaretheoremstyle[headfont=\bfseries\color{RawSienna},
mdframed={style=mdredbox},headpunct={\\[3pt]},postheadspace=0pt,]{thmredbox}

\mdfdefinestyle{mdgreenbox}{linecolor=ForestGreen,backgroundcolor=ForestGreen!5,
linewidth=2pt,rightline=false,leftline=true,topline=false,bottomline=false,}
\declaretheoremstyle[headfont=\bfseries\sffamily\color{ForestGreen!70!black},
mdframed={style=mdgreenbox},headpunct={ --- },]{thmgreenbox}

\mdfdefinestyle{mdorangebox}{linecolor=RedOrange,backgroundcolor=RedOrange!5,
linewidth=2pt,rightline=false,leftline=true,topline=false,bottomline=false,}
\declaretheoremstyle[headfont=\bfseries\sffamily\color{RedOrange!70!black},
mdframed={style=mdorangebox},headpunct={ --- },]{thmorangebox}

\mdfdefinestyle{mdblackbox}{linecolor=black,backgroundcolor=RedViolet!5!gray!5,
linewidth=3pt,rightline=false,leftline=true,topline=false,bottomline=false,}
\declaretheoremstyle[mdframed={style=mdblackbox}]{thmblackbox}

\declaretheoremstyle[spaceabove=6pt,spacebelow=6pt]{boldhead}
\declaretheoremstyle[spaceabove=6pt,spacebelow=6pt,headfont=\normalfont\itshape]{italichead}

\declaretheorem[style=boldhead,name=Problem]{problem}
\declaretheorem[style=boldhead,name=Problem,numbered=no]{problem*}
\declaretheorem[style=boldhead,name=Setup]{setup} % for config geo
\declaretheorem[style=boldhead,name=Example,sibling=problem]{example}
\declaretheorem[style=boldhead,name=$\bigstar$ Problem,sibling=problem]{niceprob}
\declaretheorem[style=boldhead,name=Theorem,sibling=problem]{theorem}
\declaretheorem[style=boldhead,name=Corollary,sibling=theorem]{corollary}
\declaretheorem[style=boldhead,name=Proposition,sibling=theorem]{prop}
\declaretheorem[style=boldhead,name=Lemma,numberwithin=problem]{lemma}
\declaretheorem[style=boldhead,name=Theorem,numbered=no]{theorem*}
\declaretheorem[style=boldhead,name=Lemma,numbered=no]{lemma*}
\declaretheorem[style=boldhead,name=Claim,numberwithin=problem]{claim}
\declaretheorem[style=boldhead,name=Claim,numbered=no]{claim*}
\declaretheorem[style=boldhead,name=Remark,numberwithin=problem]{remark}
\declaretheorem[style=boldhead,name=Remark,numbered=no]{remark*}
\declaretheorem[style=boldhead,name=Definition,sibling=theorem]{definition}
\declaretheorem[style=boldhead,name=Definition,numbered=no]{definition*}

\providecommand{\ol}{\overline}
\providecommand{\tensor}{\otimes}
\renewcommand{\hom}{\operatorname{Hom}}
\providecommand{\tor}{\operatorname{Tor}}
\providecommand{\Gal}{\operatorname{Gal}}
\providecommand{\ceil}[1]{\left\lceil #1 \right\rceil}
\providecommand{\floor}[1]{\left\lfloor #1 \right\rfloor}
\newcommand{\dd}{\,\mathrm{d}}
\providecommand{\ul}{\underline}
\providecommand{\eps}{\varepsilon}
\providecommand{\half}{\frac{1}{2}}
\providecommand{\CC}{\mathbb C}
\providecommand{\FF}{\mathbb F}
\providecommand{\II}{\mathbb I}
\providecommand{\NN}{\mathbb N}
\providecommand{\QQ}{\mathbb Q}
\providecommand{\RR}{\mathbb R}
\providecommand{\ZZ}{\mathbb Z}
\providecommand{\ind}{\mathbbm 1}
\providecommand{\dg}{^\circ}
\providecommand{\ii}{\item}
\providecommand{\setdiff}{\smallsetminus}
\providecommand{\alert}[1]{{\sffamily\textbf{\textcolor{blue}{#1}}}}
\providecommand{\inv}{^{-1}}
\providecommand{\mb}[1]{\mathbf{#1}}
\newcommand{\what}{\widehat}

\newenvironment{soln}{\begin{proof}[Solution]}{\end{proof}}

\providecommand{\pow}{\operatorname{Pow}}
\providecommand{\vocab}[1]{{\textbf{\textcolor{ForestGreen}{#1}}}}
\providecommand{\lcm}{\operatorname{lcm}}
\providecommand{\abs}[1]{\left\lvert #1\right\rvert}
\providecommand{\smabs}[1]{\lvert #1\rvert}
\providecommand{\innerprod}[1]{\langle\!\langle #1\rangle\!\rangle}
\providecommand{\norm}[1]{\left\| #1\right\|}
\providecommand{\smnorm}[1]{\| #1\|}
\providecommand{\gr}{\operatorname{gr}}

\usepackage{mathrsfs}

% place title at top right
\setlength{\droptitle}{-8ex}
\pretitle{\begin{flushright}\Large\bfseries}
\posttitle{\par\end{flushright}}
\preauthor{\begin{flushright}}
\postauthor{\end{flushright}}
\predate{\begin{flushright}}
\postdate{\end{flushright}}

\title{MAT 535 HW08}
\author{\large Aditya Pahuja}
\date{\large Due 04/16}
\begin{document}
\maketitle

\begin{problem}
    [\S13.3\#4]
    The construction of the regular heptagon amounts to
    the constructibility of $\alpha = \cos(2\pi/7)$.
    We shall see later that $\alpha$ satisfies
    the equation $x^3 + x^2 - 2x + 1 = 0$.
    Use this to prove that the regular heptagon is not constructible
    by straightedge and compass.
\end{problem}

\begin{soln}
    The polynomial $x^3 + x^2 - 2x + 1$ is irreducible
    because it has no rational roots
    (the only candidate rational roots are $\pm 1$
    by the rational root theorem, and they aren't roots).
    This means $[\QQ(\alpha) : \QQ] = 3$ is not a power of $2$,
    hence $\alpha$ is not constructible by straightedge and compass
    by Proposition 23.
\end{soln}

\begin{problem}
    [\S13.3\#5]
    Use the fact that $\alpha = 2\cos(2\pi/5)$
    satisfies the equation $x^2 + x - 1 = 0$ to conclude that
    the regular pentagon is constructible by straightedge and compass.
\end{problem}

\begin{soln}
    Since $\QQ(\alpha)$ is a quadratic extension of $\QQ$,
    $\alpha$ is constructible, and then $\sin(2\pi/5)$
    can also be constructed (since
    $[\QQ(\sqrt{1-\alpha^2}):\QQ(\alpha)]\le 2$).
    Then, we can rotate segments by $2\pi/5$,
    so we can also rotate segments by $\pi - 2\pi/5 = 3\pi/5$,
    which is the measure of the interior angles of a regular pentagon.
    Draw a segment $A_0A_1$ and, for $i\ge 0$, let $A_{i+2}$
    be the rotation $A_i$ about $A_{i+1}$ by $3\pi/5$ counterclockwise;
    this constructs the regular pentagon $A_0A_1A_2A_3A_4$.
\end{soln}

\begin{problem}
    [\S14.5\#2]
    Determine the subfields of $\QQ(\zeta_8)$ generated by
    the periods of $\zeta_8$ and in particular show that
    not every subfield has such a period as a primitive element.
\end{problem}

\begin{soln}
    The Galois group
    $G = \Gal(\QQ(\zeta_8)/\QQ)\cong(\ZZ/2\ZZ)\times(\ZZ/2\ZZ)$
    has three nontrivial subgroups $\{1,\sigma_3\}$,
    $\{1,\sigma_5\}$, and $\{1,\sigma_7\}$,
    where $\sigma_a\in G$ is the automorphism sending
    $\zeta_8$ to $\zeta_8^a$. Thus the periods of $\zeta_8$ are
    \begin{equation*}
        \zeta_8 + \zeta_8^3 = i\sqrt 2,\quad
	\zeta_8 + \zeta_8^5 = 0,\quad
	\zeta_8 + \zeta_8^7 = \sqrt 2.
    \end{equation*}
    The corresponding subfields are respectively
    $\QQ(i\sqrt 2)$, $\QQ$, and $\QQ(\sqrt 2)$,
    so the subfield $\QQ(i)$ does not have a period of $\zeta_8$
    as a primitive element.
\end{soln}

\begin{problem}
    [\S14.5\#5]
    Let $p$ be a prime and let $S_n$ be the sum
    of the $n$th powers of all the primitive $p$th roots of unity.
    Show that $S_n = -1$ when $p$ does not divide $n$
    and $S_n = p-1$ when $p$ does divide $n$.
\end{problem}

\begin{soln}
    Let $\zeta$ be a primitive $p$th root of unity, so that
    $S_n = \zeta^n + \zeta^{2n} + \cdots + \zeta^{(p-1)n}$.
    If $p$ divides $n$ then $\zeta^n = 1$, so of course $S_n = 1$.
    If $p$ does not divide $n$, then
    multiplication by $\ol n$ (the residue class of $n$)
    is an isomorphism on $(\ZZ/p\ZZ)^\times$, so
    \begin{equation*}
	S_n = \zeta^n + \zeta^{2n} + \cdots + \zeta^{(p-1)n}
	= \zeta + \zeta^2 + \cdots + \zeta^{p-1} = S_1
    \end{equation*}
    because the exponents on the left are a permutation
    of the exponents on the right (modulo $p$).
    We can then see that $S_1 = -1$ by Vieta's formulas applied to
    $\Phi_p(x) = x^{p-1} + x^{p-2} + \cdots + 1$.
\end{soln}

\begin{problem}
    [\S14.5\#7]
    Show that complex conjugation restricts to the automorphism
    $\sigma_{-1}\in\Gal(\QQ(\zeta_n)/\QQ)$ of
    the cyclotomic field of $n$th roots of unity.
    Show that the field $K^+ = \QQ(\zeta_n + \zeta_n\inv)$
    is the subfield of real elements of $K = \QQ(\zeta_n)$,
    called the \vocab{maximal real subfield} of $K$.
\end{problem}

\begin{soln}
    The conjugate of $\zeta_n$ is $\zeta_n\inv$,
    and the only automorphism of $\Gal(K/\QQ)$ sending
    $\zeta_n$ to $\zeta_n\inv$ is $\sigma_{-1}$,
    so the conjugation automorphism must restrict to $\sigma_{-1}$.
    The maximal real subfield $L$ the fixed field
    of the order-$2$ subgroup $\{1,\sigma_{-1}\}$, hence $[K:L] = 2$.
    Since $\zeta_n + \zeta_n\inv$ is real, $K^+\subseteq L$,
    hence $[K:K^+]\ge 2$. Letting $\eta = \zeta_n + \zeta_n\inv$
    and noting that $\zeta_n\zeta_n\inv = 1$,
    we can see that $\zeta_n$ is a root of the quadratic polynomial,
    \begin{equation*}
        p(t) = t^2 - \eta t + 1
    \end{equation*}
    so $[\QQ(\zeta_n):K^+] = [K:K^+]\le 2$,
    hence $[K:K^+] = [K:L] = 2$ and $K^+\subseteq L$,
    implying $K^+ = L$.
\end{soln}










\end{document}
